\section{Complete Analysis over One Hop}
\label{sec:one-hop}
The analyzer maintains one ordered work sequence and one publication boundary
per frame. On the engine thread it schedules windowing and packing,
butterflies, positive-bin reconstruction, magnitude processing, and per-bin
output. Every stage preserves the dependencies of uninterrupted evaluation, so
pausing changes where work runs, not which signal frame it analyzes. Unused
conjugate bins are never computed. One implementation,
\code{SpectrumAnalysis<T>}, supports both scalar samples and independent SIMD
channels under the same scheduling contract.

\subsection{Dependency order and work bound}
Let $M=N/2$, $K=M+1$, and $B=(M/2)\log_2 M$. Table~\ref{tab:units} defines
the operations in each stage. Let $p=4$ while window coefficients are being
rebuilt and $p=1$ when the window cache is clean, and let $o$ be a
compile-time output weight: one by default, or two for Fourier's four-channel
coordinate mapping. The ordered sequence then consumes $W$ credits, where
\begin{equation}
 W=pM+B+(1+o)K=\frac{N}{4}\log_2(N/2)+\frac{(p+1+o)N}{2}+1+o.
 \label{eq:wholework}
\end{equation}
All packing completes before any butterfly, all butterflies before any
reconstruction, and the complete prefix sum before any range smoothing, though
a single sample's quota may cross stage boundaries. Because pausing leaves the
dependency graph unchanged, the induction argument for the butterfly traversal
in Appendix~\ref{sec:legacy-transform} carries over directly. This
preservation of arithmetic order holds across pauses within one
implementation; it does not imply bitwise identity with other FFT kernels or
compiler configurations.

\begin{table}[tb]
\centering\small
\begin{tabular}{@{}rlp{0.64\linewidth}@{}}
\toprule
Credits & Stage & Scheduled work \\
\midrule
$pM$ & Pack & Each pair occupies $p$ credits: read two retained samples, prepare/apply window weights and store the permuted pair at the first credit; skip the remaining credit. \\
$B$ & Transform & Execute one complex butterfly and advance traversal indices. \\
$K$ & Reconstruct & Recover one bin, take its magnitude, and append a prefix sum if frequency smoothing is enabled. \\
$oK$ & Output & Each bin occupies $o$ credits: apply enabled smoothing and invoke the output callback at the first credit; skip the remaining credit. \\
\bottomrule
\end{tabular}
\caption{Production scheduling credits. Dirty window preparation has weight four; Fourier's coordinate output has weight two. Other operations have weight one. Credits account for ordered work, not equal-duration operations.}
\label{tab:units}
\end{table}


Call $s\in\{1,\ldots,H\}$ of a frame receives the quota
\begin{equation}
 q_s=\floor{sW/H}-\floor{(s-1)W/H}.
 \label{eq:balanced}
\end{equation}
After $s$ calls the cumulative count is $\floor{sW/H}$, so the credits are
exhausted, and the frame is published, on call $H$, with at most
$\ceil{W/H}$ credits on any call. For $H=\rho N$ with fixed $\rho>0$, this
maximum is $O(\log N)$. A weighted operation executes at its first credit and
its remaining credits are bookkeeping, so the arithmetic may finish before
publication. A quotient/remainder accumulator implements the rule without
per-sample division. The quota itself is an elementary balanced schedule; what
matters here is its application to the complete dependency-ordered analysis.
Appendix~\ref{app:one-hop} gives the dispatch listing, the proof, and examples
for clean and dirty caches.

The same telescoping argument bounds any run of calls. For fixed $W$ and $H$,
any $D$ consecutive active sample calls consume between $\floor{DW/H}$ and
$\ceil{DW/H}$ credits, including runs that cross frame boundaries. This
connects sample-level accounting to a host block and holds per analyzer,
provided the configuration is fixed and no reset or pause intervenes; it is a
bound on credits, not on duration. The logical load over any such interval
therefore differs from its proportional share by less than one credit. For a
clean scalar $N=4096$, $H=1024$ frame, $W=17410$: at most 18 credits per
sample and 1089 per 64-sample block. No allocation of $W$ credits to $H$ calls
has a smaller maximum than $\ceil{W/H}$, but that optimality concerns credits;
it says nothing about equalizing operation durations.

\begin{proposition}[Conditional operation-cost bound]
Suppose every indivisible scheduled operation $i$ has an upper cost $c_i$,
is charged integer weight $w_i\ge1$, and executes at its first credit.
Let $\kappa=\max_i c_i/w_i$ and $w_{\max}=\max_i w_i$. If capture, control,
and dispatch overhead over $D$ calls is bounded by $D c_{\mathrm{base}}(N)$,
then, under fixed configuration and without external interference, the time
$T_D$ spent in those calls satisfies
\[
 T_D\le D c_{\mathrm{base}}(N)+
 \kappa\bigl(\ceil{DW/H}+w_{\max}-1\bigr).
\]
\end{proposition}
\begin{proof}
The calls consume at most $\ceil{DW/H}$ credits. An operation started on the
last consumed credit can extend its charge by at most $w_{\max}-1$ credits;
any residual charge from an operation begun earlier only reduces new work.
Multiplying the total charged weight of newly executed operations by
$\kappa$ and adding the stated overhead proves the result.
\end{proof}
The boundary term arises because an operation executes at its first credit
rather than after all its credits have been consumed. Weights proportional to
conservative operation costs would make this a useful design rule, but
measured averages on a general-purpose operating system cannot supply such
upper bounds. The experiment measures the existing weights; it neither
calibrates $\kappa$ nor establishes a wall-time bound. A native hybrid follows
the same accounting, but its indivisible FFT call may dominate $\kappa$.

\begin{figure}[htbp]
\centering
\begin{tikzpicture}[x=0.9cm,y=0.7cm,font=\small]
\draw[->] (0,0) -- (13,0) node[right]{sample};
\draw[dashed] (1,0) -- (1,3.5);
\draw[dashed] (12,0) -- (12,3.5);
\node[below] at (1,0) {$t_r$};
\node[below] at (12,0) {$t_r+H-1$};
\node[left] at (.8,2.8) {batch};
\fill[accent!65] (1,2.4) rectangle (1.8,3.2);
\node[right] at (2,2.8) {whole pipeline; immediate logical publication};
\node[left] at (.8,1.5) {spread};
\fill[ink!25] (1,1.1) rectangle (3,1.9);
\fill[ink!50] (3,1.1) rectangle (7.8,1.9);
\fill[ink!75] (7.8,1.1) rectangle (9.5,1.9);
\fill[accent!40] (9.5,1.1) rectangle (12,1.9);
\node at (2,1.5) {pack};
\node at (5.4,1.5) {FFT};
\node[font=\scriptsize,white] at (8.65,1.5) {magnitudes};
\node[font=\scriptsize] at (10.75,1.5) {smooth/output};
\draw[->,thick] (12,1.1) -- (12,.15);
\end{tikzpicture}
\caption{One selected frame, two placements of its complete analysis. The
spread path publishes only after the final logical credit. Stage widths are
schematic; credit balance does not imply equal execution time. Native hybrids
retain one indivisible FFT inside the distributed sequence.}
\label{fig:whole-timeline}
\end{figure}


The per-call bound covers scheduled window and cache rebuilds and the modules'
bounded output callbacks. Construction and host lifecycle callbacks are
excluded, and each call still performs constant input and control work plus an
$O(\log N)$ plan-index calculation. Because every stage is decomposed into
bounded operations, no $O(N)$ boundary pass remains in steady-state engine
analysis; the quota decides only where those operations run
(Figure~\ref{fig:whole-timeline}). Reweighting changes that placement, but it
cannot equalize operation costs, constrain an arbitrary user callback, or
bound operating-system interruptions.

\subsection{Retained samples, plans, and publication time}
Count continuously captured samples from zero and let the frame endpoint be
$t_r=rH$. The frame is selected on call $t_r$; its last credit and its
publication occur on call $t_r+H-1$. The last actual output may precede the
final credit, but it remains private until publication. The nominal
sample-index ages at publication are therefore
\begin{equation}
 A_{\mathrm{newest}}=(H-1)/f_s,\qquad
 A_{\mathrm{mid}}=((N-1)/2+H-1)/f_s.
 \label{eq:wholeage}
\end{equation}
These ages exclude elapsed computation and display delay. Successive
publications are exactly $H$ samples apart, in contrast to the original
scheduler's next-call publication in Equation~\eqref{eq:age}. At startup the
window is zero-padded rather than waiting for a full window of input. When the
hop changes, endpoints satisfy $t_{r+1}=t_r+H_r$, and each frame uses its
latched $H_r$.

Because packing is incremental, the selected samples must survive while
capture continues. A preallocated ring of capacity $C=N_{\max}+H_{\max}$
suffices for the prepared maximum window and hop, and no whole-frame copy of
the input is needed. Appendices~\ref{app:one-hop} and~\ref{app:implementation}
give the lifetime proof and the reuse of maximum-size tables.

\subsection{Smoothing, configuration, and output ownership}
Reconstruction computes magnitudes and, when frequency smoothing is enabled,
a prefix sum that is complete before any range average is emitted. Output then
applies an inclusive fractional-octave magnitude mean, optional temporal
averaging, and the module's display conversion. These are magnitude means, not
power means or a constant-Q analysis, and the window normalization and display
units are unchanged. Appendix~\ref{app:one-hop} records the equations, control
conventions, cache behavior, and the sample-level quantization of the hop.

Settings latch at the start of each frame, and a reset or sample-rate change
cancels any partial analysis. Window and smoothing caches rebuild inside
scheduled units, and a length change logically clears the input and averaging
history. Storage is prepared outside the sample path, with one exception:
increasing the supported hop in a sample-rate callback can allocate. Freeze
behaves differently in the two modules, so the continuous-input timestamp
formulas do not hold across a pause.

Each output bin is written to producer-owned storage. Only a complete frame is
published, through a single-producer/single-consumer triple-slot mailbox that
uses acquire/release atomic exchange. The consumer keeps its slot until its
next successful exchange, and the producer never writes to that slot, so a
slow consumer may miss intermediate snapshots but never reads a partially
written one. Fourier uses one curve mailbox; Spectre uses one mailbox per
history column and builds and recolors its image history on the UI thread.
Lifecycle work, other shared controls, UI work, and host scheduling lie
outside the quota, so the contract is not a general claim of hard real-time
safety.
